% =====================================================================
% 8. Design and Architecture
% =====================================================================
\section{Design and Architecture}
\label{sec:design}

The system comprises three layers, with data flowing exclusively from the rig towards the browser (Figure~\ref{fig:architecture}).
No component is able to issue commands to the rig, and NFR1 is therefore satisfied by the structure of the system itself.

\begin{figure}[H]
    \centering
    \resizebox{\textwidth}{!}{% High-level architecture: three layers, adapters, one-way data flow.
\begin{tikzpicture}[
  font=\small,
  box/.style={draw=accent, fill=accent!8, rounded corners=2pt, align=center, minimum height=0.9cm, text width=2.6cm},
  ext/.style={draw=black!50, fill=black!5, rounded corners=2pt, align=center, minimum height=0.9cm, text width=2.6cm},
  layer/.style={draw=black!25, dashed, rounded corners=4pt, inner sep=6pt},
  lbl/.style={font=\footnotesize\bfseries\color{black!60}, anchor=south west},
  flow/.style={-{Stealth[length=2.2mm]}, thick, accent},
  node distance=0.45cm and 0.6cm
]
% Sources
\node[ext] (rig) {STT test rig\\\footnotesize(read-only)};
\node[ext, below=of rig] (sim) {BMS simulator};
\node[ext, below=of sim] (rep) {Recorded traces\\\footnotesize(files)};

% Ingestion
\node[box, right=1.2cm of rig] (a1) {Live-rig adapter};
\node[box, right=1.2cm of sim] (a2) {Simulator adapter};
\node[box, right=1.2cm of rep] (a3) {Replay adapter};

% Processing
\node[box, right=1.2cm of a1] (cal) {Calibration\\\footnotesize(STT logic)};
\node[box, below=of cal] (cond) {Signal conditioning};
\node[box, below=of cond] (kpi) {KPI engine};
\node[box, below=of kpi] (alert) {Alert engine\\\footnotesize(4 tiers + rules)};
\node[box, below=0.45cm of a3] (store) {Time-series store};

% Presentation
\node[box, right=1.2cm of cond] (ws) {WebSocket\\gateway};
\node[box, above=of ws] (eng) {Engineer view};
\node[box, below=of ws] (lead) {Project-lead view};
\node[box, below=of lead] (demo) {Demo view};

% Layers
\begin{scope}[on background layer]
\node[layer, fit=(a1)(a3)] (L1) {};
\node[layer, fit=(cal)(alert)] (L2) {};
\node[layer, fit=(eng)(demo)(ws)] (L3) {};
\end{scope}
\node[lbl] at (L1.north west) {Ingestion};
\node[lbl] at (L2.north west) {Processing};
\node[lbl] at (L3.north west) {Presentation};

% Flows
\draw[flow] (rig) -- (a1);
\draw[flow] (sim) -- (a2);
\draw[flow] (rep) -- (a3);
\draw[flow] (a1.east) -- (cal.west);
\draw[flow] (a2.east) -| ([xshift=-0.3cm]cal.west) -- (cal.west);
\draw[flow] (a3.east) -| ([xshift=-0.3cm]cal.west);
\draw[flow] (cal) -- (cond);
\draw[flow] (cond) -- (kpi);
\draw[flow] (kpi) -- (alert);
\draw[flow] (kpi.east) -- ++(0.35,0) |- (ws.west);
\draw[flow] (alert.east) -- ++(0.35,0) |- (ws.west);
\draw[flow] (ws) -- (eng);
\draw[flow] (ws) -- (lead);
\draw[flow] (ws.east) -- ++(0.25,0) |- (demo.east);
\draw[flow, black!50] (kpi.west) -- ++(-0.5,0) |- (store.east);
\end{tikzpicture}
}
    \caption{High-level architecture. Arrows indicate the only direction in which data can flow.}
    \label{fig:architecture}
\end{figure}

The \textbf{ingestion layer} encapsulates each data source in an adapter that implements a common interface (Listing~\ref{lst:adapter}).
Because the simulator, replay files and live rig are indistinguishable to the remainder of the system, every experiment can be conducted without access to the rig, and rig-specific code is confined to a single location.

In the \textbf{processing layer}, raw samples are passed through the STT's calibration code and then through the selected filter.
KPIs are subsequently calculated, and the alert engine determines whether an alert should be raised.
Both raw and filtered values are retained.

The \textbf{presentation layer} transmits downsampled KPI frames and alert events to the browser over WebSockets.
The three role-specific views are composed from a shared component library, each with different default settings for detail and alert visibility.

\begin{lstlisting}[language=Python, caption={Planned adapter interface. The absence of write methods ensures that the dashboard cannot affect the rig.}, label=lst:adapter]
class SourceAdapter(Protocol):
    def describe(self) -> SourceInfo: ...   # channels, rates, units
    async def stream(self) -> AsyncIterator[Sample]: ...
\end{lstlisting}

\begin{table}[H]
    \centering\small
    \caption{Principal design decisions and trade-offs.}
    \label{tab:decisions}
    \begin{tabularx}{\textwidth}{lY}
        \toprule
        Decision & Trade-off \\
        \midrule
        Read-only adapters, not a two-way twin & No remote control, but no risk to the rig \\
        Reuse of STT calibration & Dependence on partner code, but consistent values \\
        Retention of raw and filtered data & Greater storage, but faults cannot be concealed \\
        Server-side downsampling & Greater server load, but far less browser load \\
        \bottomrule
    \end{tabularx}
\end{table}

The live-rig interface will be agreed with the STT; owing to the adapter layer, this does not delay other work.
The database and charting library will be confirmed through the Experiment~1 benchmarks.
